\documentclass[12pt]{article}

% Separate working manuscript. Canonical main.* files are unchanged.
% !TEX TS-program = xelatex
% File: preamble.tex
% Purpose: Brett Reynolds house style LaTeX preamble
% Version: 2.0.0 (December 2025 typography redesign)
%
% Usage: % !TEX TS-program = xelatex
% File: preamble.tex
% Purpose: Brett Reynolds house style LaTeX preamble
% Version: 2.0.0 (December 2025 typography redesign)
%
% Usage: % !TEX TS-program = xelatex
% File: preamble.tex
% Purpose: Brett Reynolds house style LaTeX preamble
% Version: 2.0.0 (December 2025 typography redesign)
%
% Usage: \input{.house-style/preamble.tex} in main document
%
% Compiler: XeLaTeX (default) or LuaLaTeX
%           NOT pdfLaTeX - requires fontspec for fonts

% ===========================
% FONTS
% ===========================
\usepackage{fontspec}

\setmainfont{EBGaramond}[
  Extension=.otf,
  UprightFont=*-Regular,
  ItalicFont=*-Italic,
  BoldFont=*-Bold,
  BoldItalicFont=*-BoldItalic,
  Numbers=OldStyle,
  Ligatures=TeX,
]

% IPA fallback font
\newfontfamily\ipafont{CharisSIL}[
  Extension=.ttf,
  UprightFont=*-Regular,
  ItalicFont=*-Italic,
  BoldFont=*-Bold,
  BoldItalicFont=*-BoldItalic,
]
\newcommand{\ipa}[1]{{\ipafont #1}}

% Lining figures when needed (tables, years in isolation)
\providecommand{\liningnums}[1]{{\addfontfeatures{Numbers=Lining}#1}}

% Monospace
\setmonofont{Inconsolatazi4}[
  Extension=.otf,
  UprightFont=*-Regular,
  BoldFont=*-Bold,
  Scale=MatchLowercase,
]

% ===========================
% PAGE LAYOUT
% ===========================
\usepackage[
  letterpaper,
  inner=1.25in,
  outer=1in,
  top=1in,
  bottom=1.25in,
  marginparwidth=0.6in,
]{geometry}

\usepackage[british]{babel}
\usepackage[final]{microtype}

% ===========================
% HEADINGS
% ===========================
\usepackage{titlesec}

% Section: small caps, number in margin
\titleformat{\section}
  {\normalfont\scshape}
  {\llap{\thesection\quad}}
  {0pt}
  {}

% Subsection: small caps sentence case
\titleformat{\subsection}
  {\normalfont\scshape}
  {\thesubsection\quad}
  {0pt}
  {}

% Subsubsection: upright
\titleformat{\subsubsection}
  {\normalfont}
  {\thesubsubsection\quad}
  {0pt}
  {}

% Spacing
\titlespacing*{\section}{0pt}{2ex plus 1ex minus .2ex}{1ex plus .2ex}
\titlespacing*{\subsection}{0pt}{1.5ex plus 1ex minus .2ex}{0.5ex plus .2ex}
\titlespacing*{\subsubsection}{0pt}{1ex plus 0.5ex minus .1ex}{0.3ex plus .1ex}

% ===========================
% RUNNING HEADS
% ===========================
\usepackage{fancyhdr}
\pagestyle{fancy}
\fancyhf{}
\setlength{\headheight}{13.6pt}
\renewcommand{\headrulewidth}{0pt}

% For oneside documents (most papers):
\fancyhead[L]{\small\scshape\leftmark}
\fancyhead[R]{\small\thepage}
\fancyfoot{}

% ===========================
% COLORS & HYPERLINKS
% ===========================
\usepackage{xcolor}
\definecolor{linkmaroon}{RGB}{128, 0, 32}

\usepackage{hyperref}
\hypersetup{
  colorlinks=true,
  linkcolor=linkmaroon,
  citecolor=linkmaroon,
  urlcolor=linkmaroon,
  pdfauthor={Brett Reynolds},
}

\usepackage{orcidlink}

% ===========================
% QUOTATIONS
% ===========================
\usepackage[style=american]{csquotes} % Double quotes outer (Oxford spelling kept via babel)
\MakeOuterQuote{"}

% ===========================
% SEMANTIC MACROS
% ===========================
% Terms (linguistic concepts being introduced/defined)
\newcommand{\term}[1]{\textsc{#1}}

% Mentions (words/expressions under discussion)
\newcommand{\mention}[1]{\textit{#1}}

% Mentions in headings (angle brackets instead of italics)
% Robust: headings are moving arguments (written to .toc/.aux)
\DeclareRobustCommand{\mentionhead}[1]{⟨\textup{#1}⟩}

% Object language (cited forms, foreign words)
\newcommand{\olang}[1]{\textit{#1}}

% Small-caps abbreviations for glosses
\newcommand{\abbr}[1]{\textsc{#1}}

% Cross-linguistic subscript marker
\usepackage{marvosym}
\newcommand{\crossmark}{\textsubscript{\Cross}}

% ===========================
% MATHS AND SYMBOLS
% ===========================
\usepackage{amsmath,amssymb}

% ===========================
% LINGUISTIC EXAMPLES
% ===========================
\usepackage{langsci-gb4e}
\makeatletter
\@ifundefined{noautomath}{}{\noautomath}
\makeatother

% Judgement markers
\newcommand{\ungram}[1]{*\!#1}
\newcommand{\marg}[1]{?\!#1}
\newcommand{\odd}[1]{\#\!#1}

% ===========================
% LISTS
% ===========================
\usepackage{enumitem}
\setlist{nosep, leftmargin=*}
\setlist[enumerate]{label=\arabic*.}
\setlist[itemize]{label=--}

% ===========================
% TABLES & FIGURES
% ===========================
\usepackage{booktabs}
\usepackage{array}
\usepackage{graphicx}
\graphicspath{{figures/}}

% ===========================
% BIBLIOGRAPHY
% ===========================
\usepackage[backend=biber,style=apa,natbib=true,doi=true,isbn=false,url=true]{biblatex}
\addbibresource{references.bib}
\IfFileExists{references-local.bib}{\addbibresource{references-local.bib}}{}

% Possessive citation: Author's (Year)
\newcommand{\posscite}[1]{\citeauthor{#1}'s (\citeyear{#1})}

% ===========================
% UTILITIES
% ===========================
\usepackage{xspace}
\newcommand{\eg}{e.g.,\xspace}
\newcommand{\ie}{i.e.,\xspace}
\newcommand{\etc}{etc.\xspace}
 in main document
%
% Compiler: XeLaTeX (default) or LuaLaTeX
%           NOT pdfLaTeX - requires fontspec for fonts

% ===========================
% FONTS
% ===========================
\usepackage{fontspec}

\setmainfont{EBGaramond}[
  Extension=.otf,
  UprightFont=*-Regular,
  ItalicFont=*-Italic,
  BoldFont=*-Bold,
  BoldItalicFont=*-BoldItalic,
  Numbers=OldStyle,
  Ligatures=TeX,
]

% IPA fallback font
\newfontfamily\ipafont{CharisSIL}[
  Extension=.ttf,
  UprightFont=*-Regular,
  ItalicFont=*-Italic,
  BoldFont=*-Bold,
  BoldItalicFont=*-BoldItalic,
]
\newcommand{\ipa}[1]{{\ipafont #1}}

% Lining figures when needed (tables, years in isolation)
\providecommand{\liningnums}[1]{{\addfontfeatures{Numbers=Lining}#1}}

% Monospace
\setmonofont{Inconsolatazi4}[
  Extension=.otf,
  UprightFont=*-Regular,
  BoldFont=*-Bold,
  Scale=MatchLowercase,
]

% ===========================
% PAGE LAYOUT
% ===========================
\usepackage[
  letterpaper,
  inner=1.25in,
  outer=1in,
  top=1in,
  bottom=1.25in,
  marginparwidth=0.6in,
]{geometry}

\usepackage[british]{babel}
\usepackage[final]{microtype}

% ===========================
% HEADINGS
% ===========================
\usepackage{titlesec}

% Section: small caps, number in margin
\titleformat{\section}
  {\normalfont\scshape}
  {\llap{\thesection\quad}}
  {0pt}
  {}

% Subsection: small caps sentence case
\titleformat{\subsection}
  {\normalfont\scshape}
  {\thesubsection\quad}
  {0pt}
  {}

% Subsubsection: upright
\titleformat{\subsubsection}
  {\normalfont}
  {\thesubsubsection\quad}
  {0pt}
  {}

% Spacing
\titlespacing*{\section}{0pt}{2ex plus 1ex minus .2ex}{1ex plus .2ex}
\titlespacing*{\subsection}{0pt}{1.5ex plus 1ex minus .2ex}{0.5ex plus .2ex}
\titlespacing*{\subsubsection}{0pt}{1ex plus 0.5ex minus .1ex}{0.3ex plus .1ex}

% ===========================
% RUNNING HEADS
% ===========================
\usepackage{fancyhdr}
\pagestyle{fancy}
\fancyhf{}
\setlength{\headheight}{13.6pt}
\renewcommand{\headrulewidth}{0pt}

% For oneside documents (most papers):
\fancyhead[L]{\small\scshape\leftmark}
\fancyhead[R]{\small\thepage}
\fancyfoot{}

% ===========================
% COLORS & HYPERLINKS
% ===========================
\usepackage{xcolor}
\definecolor{linkmaroon}{RGB}{128, 0, 32}

\usepackage{hyperref}
\hypersetup{
  colorlinks=true,
  linkcolor=linkmaroon,
  citecolor=linkmaroon,
  urlcolor=linkmaroon,
  pdfauthor={Brett Reynolds},
}

\usepackage{orcidlink}

% ===========================
% QUOTATIONS
% ===========================
\usepackage[style=american]{csquotes} % Double quotes outer (Oxford spelling kept via babel)
\MakeOuterQuote{"}

% ===========================
% SEMANTIC MACROS
% ===========================
% Terms (linguistic concepts being introduced/defined)
\newcommand{\term}[1]{\textsc{#1}}

% Mentions (words/expressions under discussion)
\newcommand{\mention}[1]{\textit{#1}}

% Mentions in headings (angle brackets instead of italics)
% Robust: headings are moving arguments (written to .toc/.aux)
\DeclareRobustCommand{\mentionhead}[1]{⟨\textup{#1}⟩}

% Object language (cited forms, foreign words)
\newcommand{\olang}[1]{\textit{#1}}

% Small-caps abbreviations for glosses
\newcommand{\abbr}[1]{\textsc{#1}}

% Cross-linguistic subscript marker
\usepackage{marvosym}
\newcommand{\crossmark}{\textsubscript{\Cross}}

% ===========================
% MATHS AND SYMBOLS
% ===========================
\usepackage{amsmath,amssymb}

% ===========================
% LINGUISTIC EXAMPLES
% ===========================
\usepackage{langsci-gb4e}
\makeatletter
\@ifundefined{noautomath}{}{\noautomath}
\makeatother

% Judgement markers
\newcommand{\ungram}[1]{*\!#1}
\newcommand{\marg}[1]{?\!#1}
\newcommand{\odd}[1]{\#\!#1}

% ===========================
% LISTS
% ===========================
\usepackage{enumitem}
\setlist{nosep, leftmargin=*}
\setlist[enumerate]{label=\arabic*.}
\setlist[itemize]{label=--}

% ===========================
% TABLES & FIGURES
% ===========================
\usepackage{booktabs}
\usepackage{array}
\usepackage{graphicx}
\graphicspath{{figures/}}

% ===========================
% BIBLIOGRAPHY
% ===========================
\usepackage[backend=biber,style=apa,natbib=true,doi=true,isbn=false,url=true]{biblatex}
\addbibresource{references.bib}
\IfFileExists{references-local.bib}{\addbibresource{references-local.bib}}{}

% Possessive citation: Author's (Year)
\newcommand{\posscite}[1]{\citeauthor{#1}'s (\citeyear{#1})}

% ===========================
% UTILITIES
% ===========================
\usepackage{xspace}
\newcommand{\eg}{e.g.,\xspace}
\newcommand{\ie}{i.e.,\xspace}
\newcommand{\etc}{etc.\xspace}
 in main document
%
% Compiler: XeLaTeX (default) or LuaLaTeX
%           NOT pdfLaTeX - requires fontspec for fonts

% ===========================
% FONTS
% ===========================
\usepackage{fontspec}

\setmainfont{EBGaramond}[
  Extension=.otf,
  UprightFont=*-Regular,
  ItalicFont=*-Italic,
  BoldFont=*-Bold,
  BoldItalicFont=*-BoldItalic,
  Numbers=OldStyle,
  Ligatures=TeX,
]

% IPA fallback font
\newfontfamily\ipafont{CharisSIL}[
  Extension=.ttf,
  UprightFont=*-Regular,
  ItalicFont=*-Italic,
  BoldFont=*-Bold,
  BoldItalicFont=*-BoldItalic,
]
\newcommand{\ipa}[1]{{\ipafont #1}}

% Lining figures when needed (tables, years in isolation)
\providecommand{\liningnums}[1]{{\addfontfeatures{Numbers=Lining}#1}}

% Monospace
\setmonofont{Inconsolatazi4}[
  Extension=.otf,
  UprightFont=*-Regular,
  BoldFont=*-Bold,
  Scale=MatchLowercase,
]

% ===========================
% PAGE LAYOUT
% ===========================
\usepackage[
  letterpaper,
  inner=1.25in,
  outer=1in,
  top=1in,
  bottom=1.25in,
  marginparwidth=0.6in,
]{geometry}

\usepackage[british]{babel}
\usepackage[final]{microtype}

% ===========================
% HEADINGS
% ===========================
\usepackage{titlesec}

% Section: small caps, number in margin
\titleformat{\section}
  {\normalfont\scshape}
  {\llap{\thesection\quad}}
  {0pt}
  {}

% Subsection: small caps sentence case
\titleformat{\subsection}
  {\normalfont\scshape}
  {\thesubsection\quad}
  {0pt}
  {}

% Subsubsection: upright
\titleformat{\subsubsection}
  {\normalfont}
  {\thesubsubsection\quad}
  {0pt}
  {}

% Spacing
\titlespacing*{\section}{0pt}{2ex plus 1ex minus .2ex}{1ex plus .2ex}
\titlespacing*{\subsection}{0pt}{1.5ex plus 1ex minus .2ex}{0.5ex plus .2ex}
\titlespacing*{\subsubsection}{0pt}{1ex plus 0.5ex minus .1ex}{0.3ex plus .1ex}

% ===========================
% RUNNING HEADS
% ===========================
\usepackage{fancyhdr}
\pagestyle{fancy}
\fancyhf{}
\setlength{\headheight}{13.6pt}
\renewcommand{\headrulewidth}{0pt}

% For oneside documents (most papers):
\fancyhead[L]{\small\scshape\leftmark}
\fancyhead[R]{\small\thepage}
\fancyfoot{}

% ===========================
% COLORS & HYPERLINKS
% ===========================
\usepackage{xcolor}
\definecolor{linkmaroon}{RGB}{128, 0, 32}

\usepackage{hyperref}
\hypersetup{
  colorlinks=true,
  linkcolor=linkmaroon,
  citecolor=linkmaroon,
  urlcolor=linkmaroon,
  pdfauthor={Brett Reynolds},
}

\usepackage{orcidlink}

% ===========================
% QUOTATIONS
% ===========================
\usepackage[style=american]{csquotes} % Double quotes outer (Oxford spelling kept via babel)
\MakeOuterQuote{"}

% ===========================
% SEMANTIC MACROS
% ===========================
% Terms (linguistic concepts being introduced/defined)
\newcommand{\term}[1]{\textsc{#1}}

% Mentions (words/expressions under discussion)
\newcommand{\mention}[1]{\textit{#1}}

% Mentions in headings (angle brackets instead of italics)
% Robust: headings are moving arguments (written to .toc/.aux)
\DeclareRobustCommand{\mentionhead}[1]{⟨\textup{#1}⟩}

% Object language (cited forms, foreign words)
\newcommand{\olang}[1]{\textit{#1}}

% Small-caps abbreviations for glosses
\newcommand{\abbr}[1]{\textsc{#1}}

% Cross-linguistic subscript marker
\usepackage{marvosym}
\newcommand{\crossmark}{\textsubscript{\Cross}}

% ===========================
% MATHS AND SYMBOLS
% ===========================
\usepackage{amsmath,amssymb}

% ===========================
% LINGUISTIC EXAMPLES
% ===========================
\usepackage{langsci-gb4e}
\makeatletter
\@ifundefined{noautomath}{}{\noautomath}
\makeatother

% Judgement markers
\newcommand{\ungram}[1]{*\!#1}
\newcommand{\marg}[1]{?\!#1}
\newcommand{\odd}[1]{\#\!#1}

% ===========================
% LISTS
% ===========================
\usepackage{enumitem}
\setlist{nosep, leftmargin=*}
\setlist[enumerate]{label=\arabic*.}
\setlist[itemize]{label=--}

% ===========================
% TABLES & FIGURES
% ===========================
\usepackage{booktabs}
\usepackage{array}
\usepackage{graphicx}
\graphicspath{{figures/}}

% ===========================
% BIBLIOGRAPHY
% ===========================
\usepackage[backend=biber,style=apa,natbib=true,doi=true,isbn=false,url=true]{biblatex}
\addbibresource{references.bib}
\IfFileExists{references-local.bib}{\addbibresource{references-local.bib}}{}

% Possessive citation: Author's (Year)
\newcommand{\posscite}[1]{\citeauthor{#1}'s (\citeyear{#1})}

% ===========================
% UTILITIES
% ===========================
\usepackage{xspace}
\newcommand{\eg}{e.g.,\xspace}
\newcommand{\ie}{i.e.,\xspace}
\newcommand{\etc}{etc.\xspace}

\usepackage{xurl}
\usepackage[normalem]{ulem}
\setlength{\ULdepth}{1.8pt}
\setlength{\emergencystretch}{3em}
\setlength{\headheight}{14pt}
\newcommand{\exmention}[1]{\mention{#1}}
\newcommand{\exfocus}[1]{\uline{#1}}
\newcommand{\intj}[1]{\textsc{interjection\textsubscript{#1}}}
\newcommand{\intjshort}[1]{Intj\textsubscript{#1}}
\addbibresource{inquiry-relative-categories.bib}

\newif\ifblind
\blindfalse
\ifblind
\DeclareSourcemap{
  \maps[datatype=bibtex]{
    \map[overwrite=true]{
      \step[fieldsource=entrykey, match=\regexp{\Areynolds}, final]
      \step[fieldset=author, fieldvalue={Author}]
      \step[fieldset=title, fieldvalue={[Title removed for anonymous review]}]
      \step[fieldset=note, null]
      \step[fieldset=url, null]
      \step[fieldset=doi, null]
    }
  }
}
\fi
\hypersetup{
  pdftitle={Inquiry-Relative Categories: English Interjections as a Worked Case},
  pdfkeywords={inquiry-relative categories, category pluralism, projectibility, interjections, sign-types, path dependence},
}
\ifblind
\hypersetup{pdfauthor={}}
\else
\hypersetup{pdfauthor={Brett Reynolds}}
\fi
\title{Inquiry-Relative Categories:\\English Interjections as a Worked Case}
\ifblind
\author{[Author removed for blind review]}
\else
\author{Brett Reynolds \orcidlink{0000-0003-0073-7195}%
\thanks{Contact: \href{mailto:brett.reynolds@humber.ca}{brett.reynolds@humber.ca}}\\
Humber Polytechnic \& University of Toronto}
\fi
\date{\today}

\begin{document}
\maketitle

\begin{abstract}
Disputes about linguistic categories often combine different questions about
the same expressions. English interjections provide a worked case: grammatical
distribution, expressive meaning, and sequential use motivate overlapping
classifications whose boundaries needn't coincide. I develop an account of
inquiry-relative category standing in which each proposal specifies its
bearers, grouping basis, intended contribution, scope, and alternatives.
The three classifications illustrate an open family. Phonological organization
and historical ancestry motivate further candidates; social regulation and
processing connect the traditional interjection domain to broader classes.
Classifications of realizations and historical changes require explicit
relations to classifications of sign-types. These distinctions explain how
category disputes can be separated without making classification arbitrary
or assigning automatic priority to a discipline. Documented recruitment
histories support particular explanations, while audits of two corpus studies
withdraw their intended confirmatory interpretations. The resulting framework
distinguishes descriptive organization, projection, representational benefit,
and causal explanation as contributions requiring appropriate evidence.
Overlap, gradience, and shared mechanisms are compatible with justified
classification; the number and reach of useful categories remain outcomes
of inquiry.
\end{abstract}

\noindent\textbf{Keywords:} inquiry-relative categories; category pluralism;
projectibility; interjections; sign-types; path dependence

\aidisclosure{Claude Opus 4.8, Claude Opus 5, Claude Opus 5.5, GPT-5.6 Sol,
GPT-6 Astra Pro, and GPT-6 (Codex)}

\section{Introduction}\label{sec:intro}

Standalone \mention{Damn!}, complement-taking \mention{Damn these mosquitoes!},
and prenominal \mention{damn} in \mention{the damn car} present different
combinations of grammatical and semantic properties. The modifier participates
in a noun phrase while retaining an expressive contribution. The standalone
form can constitute a stance display. The complement-taking expression has
internal structure even when the whole expression stands apart from a
surrounding clause. A single judgment about whether \mention{damn} is an
interjection leaves these differences unresolved.

Which judgment matters depends on the inquiry. A grammarian investigating
phrase structure needs a distinction that a semanticist studying expressive
contribution may properly cross. An analyst of conversation asks a further
question about conventional actions and recipients' responses. The difficulty
arises when evidence answering one question is taken to settle all three,
or when failure to share one extension is treated as failure to identify
any defensible category.

This paper develops an account of \term{inquiry-relative classification}.
Its central claim is that a traditional linguistic domain can contain several
overlapping classifications, each assessed for the distinctions, inferences,
or explanations it supports. An inquiry identifies a question and a standard
of success. Evidence constrains which grouping answers it. The same objects
can participate in several justified classifications without one inquiry
acquiring general priority over the others.

English interjections supply the worked case. I develop a grammar-facing
\intj{syn}, a semantics-facing \intj{sem}, and an interaction-facing
\intj{prag}. These are analytical reconstructions drawing on established
traditions, rather than three technical senses already distinguished
consistently in the literature. They share conventionalized sign-types as
bearers: identifiable expressions in conventional uses. The three uses of
\mention{damn} above belong to one historical family but require separate
analysis at this grain.

The three classifications are examples within an open family. Phonological
organization can motivate another grouping of sign-types, and historical
ancestry can relate descendants with different present distributions. Other
inquiries classify realizations, developmental episodes, or processes of
change. Their connection to sign-types needs to be stated. These extensions
show why the subscript in \intj{x} identifies a particular inquiry, rather
than allocating one category to each disciplinary department.

The general proposal has two constraints. First, a candidate specifies its
bearer, grouping basis, intended contribution, scope, and relevant alternatives.
Second, the evidence supporting one contribution isn't automatically evidence
for another. A useful descriptive contrast, a successful prediction, an
economical representation, and a supported causal explanation can come apart.
The standards apply equally to the initial three classifications and to
further candidates.

The problem is familiar in part from the distinction between grammatical
clause type and speech act. An interrogative can perform a request, while
an answer-seeking action can be performed without interrogative form. The
proper-noun/proper-name distinction offers another comparison: a name can
be headed by a common noun. The interjection case develops the general
problem further because several traditions compete for the same label,
and their diagnoses mix properties of forms, uses, occurrences, and histories.

The historical argument concerns \term{path dependence}: an expression's
development can leave source-specific traces after its use converges with
that of other expressions. Gehweiler's history of \mention{gee} supplies a
documented case; the comparison with complement-taking \mention{damn} is
partly reconstructive. These histories help distinguish an explanatory
classification by ancestry from a classification by current distribution.
They also show why a present trace needs independent historical evidence
before it can count as support for the proposed route.

The corpus material has a narrower role. Geographic frequency differences
and changes in tagged inventories provide examples of what a measure can
and can't establish about a classification. Audits invalidate the intended
confirmatory interpretations of two further analyses, of derivatives and
\mention{fie}; the appendices preserve their records and explain the failures.
The paper uses those failures to delimit its evidence rather than treating
precise numerical outputs as grounds for category standing.

Section~\ref{sec:worked} develops the three classifications through concrete
contrasts. Section~\ref{sec:assessment} explains their common assessment
standard. Sections~\ref{sec:phonology} and~\ref{sec:history} develop phonological
and historical extensions, including the relations among different bearers.
Section~\ref{sec:extensions} considers broader social and processing classes.
The empirical assessment and discussion then establish the limits and
consequences of the general account.

\section{Three classifications of conventional signs}\label{sec:worked}

Grammars and studies of meaning and interaction have organized the traditional
interjection domain differently. Jespersen treated interjection as a manner
in which words of other categories are used \citep[90]{jespersen1924}.
\textit{The Cambridge Grammar of the English Language} (\textit{CGEL})
recognizes a lexical category \citep[22]{HuddlestonPullum2002}. The
\textit{Longman Grammar of Spoken and Written English} places interjections
within the broader class of inserts, while assigning \mention{damn},
\mention{gosh}, and \mention{gee} to expletives \citep[1082--1083, 1094--1095]{Biber1999}.

The disputes concern the organization of a well-studied domain. Libert
questions the common description of interjections as neglected
\citep{libert2019}; Cram identifies their recurring role in controversies
over the syntax--pragmatics and semantics--pragmatics boundaries
\citep[63]{cram2008}. A comparison of classifications should preserve the
contrasts those controversies disclose. It should also distinguish an
alternative grouping from a disagreement about the facts.

\subsection{Grammar and expressive meaning}

The grammar-facing classification begins with standalone and supplement
distribution. In \mention{Damn, we're going to be late!}, \mention{damn}
is appended to the clause without being its complement or modifier.
\textit{CGEL} treats interjections as occurring on their own or functioning
as supplements \citep[1350, 1361]{HuddlestonPullum2002}. In their conventional
interjectional uses they are characteristically invariant and resist the
ordinary constituent positions available to nouns, verbs, or adjectives.

Supplement function is shared by several categories. Compare the interjection
with the relative clause and noun phrase in the following examples, adapted
from \textit{CGEL} \citep[1356--1357, 1360]{HuddlestonPullum2002}.

\ea\label{ex:supplements}
  \ea \exmention{\exfocus{Ah}, so you were there after all!}
  \ex \exmention{We called in to see Sue's parents, \exfocus{which made us rather late}.}
  \ex \exmention{A university lecturer, \exfocus{Dr Brown}, was arrested for the crime.}
  \z
\z

Their external relation to a host doesn't erase their internal grammatical
differences. A nominal supplement remains a noun phrase; a relative clause
has clause structure. For the same reason, the external detachment of
\mention{Damn these mosquitoes!} is compatible with internal complementation.
\textit{CGEL} tentatively treats such cases as exceptional interjections
heading an interjection phrase with a noun-phrase complement
\citep[1361 n.~70]{HuddlestonPullum2002}. A verbal analysis remains a substantive
alternative. The complement alone doesn't decide between them.

Prenominal \mention{damn} provides a different contrast. In
\mention{the damn car}, it occurs inside the noun phrase, with the phrase's
prosodic packaging, while contributing an attitude toward its referent.
Potts's analysis distinguishes that expressive contribution from descriptive
content \citep[171, 182]{potts2007}. The expression can preserve this
semantic distinction while departing from the standalone distribution
characteristic of grammatical interjections.

\intj{sem} develops this expressive classification. \term{Expressive} here
concerns a conventional contribution distinguishable from descriptive content;
the associated \term{use-conditional} analysis concerns what using the
expression conveys. Potts identifies several diagnostics, including
perspective dependence, immediacy, and intensifying repetition
\citep[166--167]{potts2007}. Their applicability to a proposed member is an
empirical question. Mere independence from the host proposition is
insufficient, since a supplement can contribute a descriptive proposition
of its own.

The two \mention{damn} uses make the benefit of separate classification
visible. The grammatical analysis preserves a difference in distribution;
the semantic analysis preserves a common kind of contribution. Averaging
these observations into a single degree of interjectionhood would discard
which respect is shared and which differs. A constructional inventory could
also record both relations. The proposal concerns their organization as
distinct classifications and the consequences attributed to them.

\subsection{Interaction and its subtypes}

An interaction-facing classification groups conventional minimal resources
by their availability for identifiable jobs in sequences of action.
Continuers such as \mention{mm-hm}, repair initiators such as \mention{huh?},
and change-of-state uses of \mention{oh} are central cases in accounts of
interjections in conversation \citep[257--259, 263--265]{dingemanse2024}.
Their importance becomes less visible if the entire domain is organized
around emotional exclamation.

A continuer can support another speaker's ongoing turn. A repair initiator
can make a response to a problem of hearing or understanding relevant.
Those are different actions with different consequences. Assigning both
to \intj{prag} preserves their availability as conventional resources;
it doesn't predict the same next turn for both. The evidential question
for the superordinate class concerns what generalization it supports beyond
the separate descriptions of its subtypes.

The selected continuer and repair uses also show why an interactional
contribution needn't be an expressive one. A speaker can convey an attitude
through a particular \mention{mm-hm} or \mention{huh?} without that attitude
being the conventional expressive contribution of the use. I provisionally
exclude the selected readings from \intj{sem}, while assigning them to
\intj{syn} on their conventional standalone distribution and to \intj{prag}
on their sequential role. Evidence for another reading could warrant another
assignment.

Routine formulae expose disagreement within this interactional inquiry.
Ameka separates forms such as \mention{bye}, \mention{sorry}, and
\mention{thank you} from interjections, emphasizing addressees, conventional
social situations, and speech-act status \citep[108--110]{ameka1992}.
Wilkins treats them as a pragmatic and semantic subtype
\citep[142]{wilkins1992}. Both attend to meaning and interaction. Allocating
one author to grammar and the other to pragmatics would misdescribe the
disagreement.

Under the criterion adopted here, closing \mention{bye} belongs to
\intj{prag}: it is a conventional minimal resource in a recognizable
sequential environment. Closing formulae remain distinguishable from
continuers, repair initiators, and stance displays. This broad assignment
is a proposal whose usefulness can be assessed against narrower groupings.
An inquiry can support several classifications, including subtypes, rather
than requiring one extension for all interactional purposes.

\subsection{The common bearer and the contrasts it preserves}

A \term{sign-type} is an identifiable expression in a conventional use.
Its \term{profile} is the combination of grammatical, semantic, interactional,
and other properties attributed to it. Profile is a descriptive term here.
The use can retain its identity while a proposed property turns out to be
absent. Otherwise every surprising observation would identify a different
object and leave the original classification immune to correction.

Different conventional uses require evidence for their individuation.
The recurrent contrast between standalone and prenominal \mention{damn}
supports distinguishing them. A spelling alone doesn't: a word family can
include several signs, and a sign can have spelling variants. Table~\ref{tab:contrasts}
records the selected uses and the present analytical assignments. They are
neither counts of occurrences nor results of a new annotation study.

\begin{table}[htbp]
\centering
\small
\caption{Selected uses under the three worked classifications. A check
marks assignment, a dash exclusion under the stated analysis, and a question
mark an unresolved assignment. The prenominal row concerns the head use of
\mention{damn}; external distribution is assessed in its phrase.}
\label{tab:contrasts}
\begin{tabular}{@{}lccc@{}}
\toprule
Selected use & \intjshort{syn} & \intjshort{sem} & \intjshort{prag} \\
\midrule
\mention{gee!} & $\checkmark$ & $\checkmark$ & $\checkmark$ \\
\mention{damn!} (stance) & $\checkmark$ & $\checkmark$ & $\checkmark$ \\
\mention{damn these mosquitoes!} & ? & $\checkmark$ & ? \\
Prenominal \mention{damn} & -- & $\checkmark$ & -- \\
\mention{mm-hm} (continuer) & $\checkmark$ & -- & $\checkmark$ \\
\mention{huh?} (repair) & $\checkmark$ & -- & $\checkmark$ \\
\mention{bye} (closing) & $\checkmark$ & -- & $\checkmark$ \\
\mention{Bang!} (depiction) & $\checkmark$ & -- & ? \\
\bottomrule
\end{tabular}
\end{table}

The depictive use of \mention{Bang!} illustrates another separation.
Standalone distribution motivates its grammatical assignment in the analysis
adopted here. Depicting a sound supplies a different semantic relation from
the expressive contribution of \mention{damn}. Work on ideophones and
onomatopoeia makes that difference consequential
\citep{dingemanse2018,dingemanseakita2017,meinard2015}. Its conventional
sequential role remains a further question. Standalone packaging by itself
doesn't answer it.

The unresolved cases matter as much as the positive assignments.
Complement-taking \mention{damn} raises a grammatical question independently
of whether it has expressive meaning. Its absence from a dataset of
backchannels would say little about its availability for stance displays.
A question mark preserves that separation. Treating missing evidence as
exclusion would make the apparent distance between classifications partly
an artefact of what investigators happened to study.

These comparisons establish a reason to keep the judgments distinct.
They leave open how much further work each grouping performs. A grammatical
category may support additional distributional expectations; the expressive
classification may extend to signs never called grammatical interjections;
interactional subtypes may be more informative than their common superordinate
class. Those possibilities require one assessment standard that applies
without privileging any of the three.

